\scanpage{267}\mspage{183}
% archive sheet number: 266
$\pi'^{\pi}\geq\pi'_1$ implique $\pi'^{\pi}\geq\pi_1$, donc on peut appliquer
la condition c) des \unsure{fcg|feg} pour $(\mathcal F^\pi)_{\pi\geq\pi_0}$ pour
conclure qu'elle reste encore vérifiée. Donc
en somme, \uline{toutes} les conditions des fcg
(resp. fcga) sont trivialement vérifiées pour
la fcg \uline{restreinte} de la première.

De $=$, le fait qu'un cardinal $\pi'\inserted{\geq\pi'_0}$ est
utile \deleted{pour} est trivialement indépendant de
quelle des deux filtrations on considère,
elle ne dépend que de la donnée du
cardinal $\pi'$ et de la sous-catégorie $\mathcal F^{\pi'}$
de $M$. Reste la majoration des cardinaux, dans le cas où on s'intéresse à
la \uline{régularité} de $\pi'$ pour les \deleted{\unsure{filtr}} fcga
envisagées. La majoration exigée est
la $=$, mais la condition \inserted{\deleted{\unsure{card}} préalable} de
validité de cette majoration, savoir
\[
\pi'^{\pi}\geq\pi'_1\mathpunct{.}\qquad\text{resp.}\quad\pi'^{\pi}\geq\pi_1
\]
n'est pas la même. Comme $\pi'_1\geq\pi_1$,
on voit que la condition \unsure{pré=alable} est
moins restrictive dans le cas de la
première filtration, que de la seconde, d'où
le ``il suffit'' dans la dernière assertion. Mais
\uline{lorsque} $\deleted{\unsure{\pi|2}}2^\pi\geq\pi'_1$, \deleted{\unsure{or}} la plus restrictive est automatique%
